%
% Introduction.tex
%
% Copyright (C) 2025 by Inspirefly.
%
% 2.305 GHz Circular Patch Antenna Documentation
%

\chapter{Introduction}

The 1P PocketQube is a $50 \times 50 \times 50$~mm cube, with mass below 250~g. This document records the design of a 2.305~GHz centered circular microstrip patch antenna for the PocketQube parent antenna board. The frequency was selected to provide an S-band downlink/uplink channel separate from UHF telemetry, with sufficient bandwidth for higher data rates and with regulatory compatibility in the 2.25--2.35~GHz region. The antenna is realized as a probe-fed circular patch on a multilayer PCB, with two orthogonal feed points intended for circular polarization through an external quadrature power divider.

\section{Why a Circular Patch}

Rectangular and circular patches share the same operating principle: a resonant cavity formed by the metal patch, the grounded dielectric substrate, and the fringing fields at the edges. However, the circular geometry was chosen for two reasons.

First, the dominant $\mathrm{TM}_{11}$ mode of a circular patch has a field distribution with azimuthal symmetry, which simplifies generation of circular polarization by feeding at two points separated by $90^{\circ}$ in space and driven in quadrature. For a rectangular patch, circular polarization requires either a nearly square patch with truncated corners or a dual-feed with unequal slot perturbations; both are more sensitive to etching tolerances at 2.3~GHz than the symmetric circular case.

Second, for a given resonant frequency the circular patch occupies less area than the circumscribed square patch and has no sharp corners where current crowding increases loss. On a 50~mm board, a radius near 19--20~mm (see Chapter~2) leaves a usable ground-plane ring for connectors and keep-out, while a $ \lambda_g/2 $ square patch with the same substrate would require a larger footprint once fringing corrections are included.

\section{Objectives}

The design objectives for the parent antenna board are:

\begin{enumerate}[noitemsep]
    \item Resonant frequency $f_0 = 2.305$~GHz, with $S_{11} < -10$~dB across at least 20~MHz.
    \item Realized gain greater than 5~dBi at boresight (zenith) and axial ratio below 3~dB for the circularly polarized configuration.
    \item Input impedance near 50~$\Omega$ at each coax feed, compatible with the Molex 732511350 U.FL-compatible connector and the external Mini-Circuits-style divider APS-02-WBM-LP-DF-CC.
    \item Full integration into the KiCAD PCB stack without post-assembly tuning stubs, using the same fabrication stack available for the rest of the spacecraft electronics (FR-408HR multilayer).
    \item Reproducibility: all geometry originates from the ANSYS HFSS parametric model and is exported via DXF to the KiCAD footprint \texttt{Inspirefly:CircularPatchAntenna}, with version control in git.
\end{enumerate}

\section{Document Structure}

Chapter~2 develops the underlying theory and describes the three tool chains used: MATLAB Antenna Toolbox for initial analysis, ANSYS HFSS (via AEDT files in \texttt{ansys/}) for full-wave finite-element verification, and KiCAD 10.0 for layout (\texttt{KiCAD/antennaBoard/}). Chapter~3 reports the numerical results from each stage, including the limits of the early rectangular-patch estimate and the current best MATLAB optimization log (\texttt{matlab/best.txt}). Chapter~4 interprets those results and records design decisions visible in git history, such as the move from a single MMCX feed to two orthogonal coax feeds and the substrate change from Rogers RO4350B to FR-408HR. Chapter~5 summarizes status and remaining verification steps before fabrication.

\nomenclature{PCB}{Printed Circuit Board}
\nomenclature{HFSS}{High Frequency Structure Simulator}
\nomenclature{TM}{Transverse Magnetic}
\nomenclature{AR}{Axial Ratio}
\nomenclature{RL}{Return Loss, $-20\log_{10}|S_{11}|$}
